% !PRINTED 349
\weberpage{349}

Equation (8) itself now already gives a fairly good approximate value for the real root

\[
x = \sqrt[8]{82} = 1{,}73471.
\]

A more accurate value, which one obtains by going two steps further, is 1,73469, which is still correct in the fifth decimal place.

This method does not give an enclosure of the sought root between two limits. The criterion for whether a sufficient degree of accuracy has been reached consists in this: that if $s_\mu$, $s_{\mu'}$ are two successive power sums of the roots used in the computation, these are related, with sufficient accuracy, as the $\mu$th to the $\mu'$th power of a quantity, or that approximately

\[
\mu \log s_{\mu'} - \mu' \log s_\mu = 0
\]

holds. One must, however, carry out this test on several successive terms, since in special cases this relation may appear to be satisfied and then cease to hold again in the later computation.

\begin{center}
\S.\ 112.

\gesperrt{Trigonometric Solution of Cubic}

\gesperrt{Equations.}
\end{center}

We now discuss a few further methods of numerical solution applicable to equations of special forms. The aim of these methods is to make the widely available tables of the trigonometric functions and of logarithms useful for the solution of equations. We turn first to the consideration of the cubic equations, which we always assume in the reduced form

\weq{1}{x^3 + a x + b = 0}

in which $a$ and $b$ are real numbers. Since the interchange of $x$ with $-x$ is equivalent to the interchange of $b$ with $-b$, we may restrict ourselves to the assumption that $b$ is negative, and there then remain three different cases to consider. We set $b = -g$ and, according as $a$ is positive or negative, $a = \pm e$. Then we have
%CONT>>
